\BeleskeID{08-s056}
% element: 08-s056:e04
% element: 08-s056:notes
Kada se završi prelazni proces uz konstantnu pobudu, struja kondenzatora postaje nula, a pošto je to ista struja kao kroz otpornik, izlaz na otporniku teži nuli. Dijagram pokazuje pozitivan skok izlaza i njegovo eksponencijalno opadanje.

U izvođenju diferencijatorskog približenja napon kondenzatora izražava se kao razlika ulaza i napona otpornika. Diferenciranjem se dobija član sa izvodom ulaza i član sa izvodom struje. Poslednji se zanemaruje samo u režimu u kome je njegov doprinos mali u odnosu na prvi. Približenje nije opis tačnog odziva na svaku proizvoljnu pobudu.
